% BEGIN VOL07-EXPANSION VII06
\section{Supplementary worked examples}
\label{sec:vii06-pedagogy}

\begin{example}[The sphere as a regular level set]\label{ex:vii06-ped-01}
For \(F(x)=\lVert x\rVert^2\), the value \(1\) is regular because \(dF_p=2p\neq0\) on \(S^{n-1}\).  The regular level-set theorem makes \(S^{n-1}\) an embedded \((n-1)\)-submanifold and gives \(T_pS^{n-1}=p^\perp\).
\end{example}

\begin{example}[Graphs are embedded submanifolds]\label{ex:vii06-ped-02}
If \(f:\mathbb R^m\to\mathbb R^n\) is smooth, its graph is the regular level set of
\[
F(x,y)=y-f(x).
\]
The derivative with respect to \(y\) is the identity, so \(dF\) has rank \(n\) everywhere.  Hence the graph is an embedded \(m\)-submanifold.
\end{example}

\begin{example}[The diagonal encodes equality]\label{ex:vii06-ped-03}
For a smooth manifold \(M\), the diagonal \(\Delta=\{(p,p):p\in M\}\subset M\times M\) is an embedded submanifold diffeomorphic to \(M\).  Its tangent space is
\[
T_{(p,p)}\Delta=\{(v,v):v\in T_pM\}.
\]
This example later turns transversality and coincidence problems into intersection questions.
\end{example}

\section{Graded supplementary exercises}
The following exercises add a deliberate progression from concrete calculation to structural proof, hypothesis testing, application, and synthesis.

\subsection*{Standard computations and constructions}

\begin{exercise}[Circle as level set]\label{exr:vii06-ped-01}
Use the regular level-set theorem to show \(S^1\subset\mathbb R^2\) is a \(1\)-submanifold.
\end{exercise}
\begin{hint}
Take \(F=x^2+y^2\).
\end{hint}
\begin{solution}
On \(F^{-1}(1)\), \(dF=2x\,dx+2y\,dy\neq0\), so \(1\) is a regular value and the level set has dimension \(2-1=1\).
\end{solution}

\begin{exercise}[Graph dimension]\label{exr:vii06-ped-02}
Find the dimension of the graph of a smooth \(f:\mathbb R^3\to\mathbb R^2\).
\end{exercise}
\begin{hint}
Use the graph parametrization \(x\mapsto(x,f(x))\).
\end{hint}
\begin{solution}
The graph is diffeomorphic to \(\mathbb R^3\), hence has dimension \(3\).
\end{solution}

\begin{exercise}[Product tangent space]\label{exr:vii06-ped-03}
Identify \(T_{(p,q)}(M\times N)\).
\end{exercise}
\begin{hint}
Use product charts.
\end{hint}
\begin{solution}
It is canonically \(T_pM\oplus T_qN\).
\end{solution}

\begin{exercise}[Diagonal tangent space]\label{exr:vii06-ped-04}
Compute \(T_{(p,p)}\Delta\subset T_pM\oplus T_pM\).
\end{exercise}
\begin{hint}
Differentiate \(p\mapsto(p,p)\).
\end{hint}
\begin{solution}
It is \(\{(v,v):v\in T_pM\}\).
\end{solution}

\begin{exercise}[Cylinder level set]\label{exr:vii06-ped-05}
Show \(x^2+y^2=1\) defines an embedded surface in \(\mathbb R^3\).
\end{exercise}
\begin{hint}
The gradient never vanishes on the level set.
\end{hint}
\begin{solution}
For \(F=x^2+y^2\), \(dF=(2x,2y,0)\neq0\) when \(F=1\); the level set has dimension \(2\).
\end{solution}

\subsection*{Proofs}

\begin{exercise}[Graph theorem]\label{exr:vii06-ped-06}
Prove the graph of a smooth map \(f:M\to N\) is an embedded submanifold of \(M\times N\).
\end{exercise}
\begin{hint}
Use local product charts or the graph embedding \(p\mapsto(p,f(p))\).
\end{hint}
\begin{solution}
In product charts the graph is the graph of a Euclidean smooth map and is straightened by \((x,y)\mapsto(x,y-f(x))\), hence is embedded.
\end{solution}

\begin{exercise}[Product of submanifolds]\label{exr:vii06-ped-07}
If \(A\subset M\) and \(B\subset N\) are embedded submanifolds, prove \(A\times B\subset M\times N\) is embedded.
\end{exercise}
\begin{hint}
Take adapted charts for \(A\) and \(B\).
\end{hint}
\begin{solution}
The product of adapted charts sends \(A\times B\) to a coordinate subspace of the product Euclidean space.
\end{solution}

\begin{exercise}[Tangent of a product]\label{exr:vii06-ped-08}
Prove the canonical isomorphism \(T_{(p,q)}(M\times N)\cong T_pM\oplus T_qN\).
\end{exercise}
\begin{hint}
Use differentials of the two projections.
\end{hint}
\begin{solution}
The map \(v\mapsto(d\pi_Mv,d\pi_Nv)\) is an isomorphism in product coordinates, where it is the standard decomposition of \(\mathbb R^{m+n}\).
\end{solution}

\begin{exercise}[Regular level dimension]\label{exr:vii06-ped-09}
Prove that a regular level set \(F^{-1}(c)\) of \(F:M^m\to N^n\) has dimension \(m-n\).
\end{exercise}
\begin{hint}
Use the submersion normal form.
\end{hint}
\begin{solution}
Local coordinates make \(F\) a projection onto \(n\) coordinates, so the level set is locally defined by fixing those \(n\) coordinates and has \(m-n\) free coordinates.
\end{solution}

\subsection*{Counterexamples and hypothesis tests}

\begin{exercise}[Every level set is a submanifold]\label{exr:vii06-ped-10}
Test the claim with \(F(x,y)=x^2-y^2\) at level \(0\).
\end{exercise}
\begin{hint}
Inspect the origin.
\end{hint}
\begin{solution}
False.  The zero set is two crossing lines and \(dF_0=0\); the origin is not locally a one-dimensional manifold.
\end{solution}

\begin{exercise}[Injective immersion is embedded]\label{exr:vii06-ped-11}
Test the claim in general.
\end{exercise}
\begin{hint}
Recall that the topology on the image also matters.
\end{hint}
\begin{solution}
False without additional hypotheses such as properness; an injective immersion can have an image whose subspace topology does not match the source topology.
\end{solution}

\begin{exercise}[Products add codimensions incorrectly]\label{exr:vii06-ped-12}
Test whether \(\operatorname{codim}(A\times B)=\max(\operatorname{codim}A,\operatorname{codim}B)\).
\end{exercise}
\begin{hint}
Compute dimensions.
\end{hint}
\begin{solution}
False.  Codimensions add: \((m+n)-(a+b)=(m-a)+(n-b)\).
\end{solution}

\subsection*{Applications and investigations}

\begin{exercise}[Constraint manifold]\label{exr:vii06-ped-13}
Explain why independent smooth constraints \(F_1=\cdots=F_k=0\) typically define a codimension-\(k\) manifold.
\end{exercise}
\begin{hint}
Package them into \(F=(F_1,\ldots,F_k)\).
\end{hint}
\begin{solution}
If \(dF\) has rank \(k\) on the common zero set, \(0\) is a regular value and the level-set theorem gives codimension \(k\).
\end{solution}

\begin{exercise}[Configuration products]\label{exr:vii06-ped-14}
Why is the configuration space of two unconstrained particles on manifolds \(M\) and \(N\) naturally \(M\times N\)?
\end{exercise}
\begin{hint}
A configuration is an ordered pair of positions.
\end{hint}
\begin{solution}
The product manifold records the two positions independently, and its tangent space splits into the two particle velocity spaces.
\end{solution}

\subsection*{Challenge problems}

\begin{exercise}[Diagonal and fixed points]\label{exr:vii06-ped-15}
For \(f:M\to M\), express fixed points as an intersection in \(M\times M\).
\end{exercise}
\begin{hint}
Use the graph and diagonal.
\end{hint}
\begin{solution}
A point \(p\) is fixed exactly when \((p,f(p))=(p,p)\), so fixed points correspond to \(\Gamma_f\cap\Delta\).
\end{solution}

\begin{exercise}[Proper injective immersion]\label{exr:vii06-ped-16}
Explain why a proper injective immersion is an embedding.
\end{exercise}
\begin{hint}
Use that a proper continuous map into a Hausdorff space is closed.
\end{hint}
\begin{solution}
Properness makes the injective immersion a homeomorphism onto its image, while the immersion condition gives the local embedded-submanifold form; together these are precisely an embedding.
\end{solution}

% END VOL07-EXPANSION VII06
